\documentclass[11pt,a4paper]{article}

\usepackage[T1]{fontenc}
\usepackage[utf8]{inputenc}
\usepackage[spanish]{babel}
\usepackage[a4paper,margin=2.8cm]{geometry}
\usepackage{microtype}
\usepackage[hidelinks]{hyperref}

\hypersetup{
  pdftitle={El Aleph},
  pdfauthor={Juan José Gómez Cadenas}
}

\setlength{\parindent}{1.25em}
\setlength{\parskip}{0.35em}
\raggedbottom

\title{El Aleph}
\author{Juan José Gómez Cadenas}
\date{}

\begin{document}

\maketitle

Una célebre canción de Joaquín Sabina aconseja no volver jamás al lugar en el que se ha sido feliz. Confieso que mientras abría el volumen I de las obras completas de Borges ---Emecé editores, 1989--- y buscaba el texto de \emph{El Aleph} ---empieza en la página 617--- temía que la frase se aplicara también a los relatos que más nos han marcado.

Es cierto que en los últimos treinta y cinco años he abierto esos tomos a menudo y releído numerosas piezas. En algunos casos, como \emph{Hombre de la esquina rosada}, \emph{La biblioteca de Babel} o \emph{La rosa de Paracelso}, he perdido la cuenta de cuántas veces lo he hecho. Pero no había vuelto a visitar \emph{El Aleph} desde hace décadas y temía que, al hacerlo, la realidad del texto no estuviera a la altura de mis exaltados recuerdos.

No fue así. De hecho, acabé haciendo lo contrario de lo que Sabina recomienda. Necesité tres lecturas consecutivas para digerir uno de los cuentos más breves y enigmáticos del maestro argentino.

Hay una primera lectura, que, en el fondo, también es la última. \emph{El Aleph} es, en último término, una historia de amor trágico:

\begin{quotation}
La candente mañana de febrero en que Beatriz Viterbo murió, después de una imperiosa agonía que no se rebajó un solo instante ni al sentimentalismo ni al miedo, noté que las carteleras de fierro de la Plaza Constitución habían renovado no sé qué aviso de cigarrillos rubios; el hecho me dolió, pues comprendí que el incesante y vasto universo ya se apartaba de ella y que ese cambio era el primero de una serie infinita.
\end{quotation}

Era el propio Borges quien recomendaba empezar una buena historia con un nacimiento, una muerte o un forastero que llega a una ciudad sin nombre. Y la muerte de Beatriz, en la primera línea, no sólo fija el tono sombrío del relato sino que refiere al lector a esa otra Beatriz, que Dante exaltó en su \emph{Comedia}. Pero hay más. El siguiente párrafo nos informa de un imperceptible cambio que contribuye a alejar el universo de la difunta. ¿Sólo de ella? Borges invoca sutilmente la realidad del cambio constante del cosmos, que a su vez refleja el aumento incesante de la entropía. Desde la muerte de Beatriz Viterbo, en 1929, esa entropía que crece empecinadamente ha ido cambiando el universo ---con un anuncio de cigarrillos, con una guerra mundial, con la llegada del hombre a la luna, con la irrupción reciente de la IA---, pero sobre todo, lo ha ido desordenando. Ha pasado casi un siglo de los hechos que se relatan en el cuento. En la Vía Láctea, se habrán apagado alrededor de cien estrellas, más de un billón si contamos todo el universo. Desde la muerte de Beatriz, cerca de mil galaxias han emitido el último fotón que jamás nos alcanzará. Borges intuye ---y plasma, con un detalle trivial--- una tragedia cósmica. Nuestro universo se expande, empujado por la energía oscura, acelerando cada vez más, enfriándose poco a poco, hacia una muerte térmica ineludible.

Sigue, tras el párrafo inicial, lo que podríamos llamar la parte burlesca del relato, protagonizada por el primo de Beatriz, el pedante Carlos Argentino Daneri, a quien Borges describe con ironía no exenta de maldad:

\begin{quotation}
Beatriz era alta, frágil, muy ligeramente inclinada; había en su andar (si el oxímoron es tolerable) una como graciosa torpeza, un principio de éxtasis; Carlos Argentino es rosado, considerable, canoso, de rasgos finos. […] Su actividad mental es continua, apasionada, versátil y del todo insignificante.
\end{quotation}

Desde la muerte de Beatriz hasta la época en que se desarrolla el relato (1941), han pasado doce años, en los que Borges repite regularmente sus visitas al pariente que no soporta, a cambio de aferrarse a lo poco que le queda de la amada:

\begin{quotation}
Beatriz Viterbo, de perfil, en colores; Beatriz, con antifaz, en los carnavales de 1921; la primera comunión de Beatriz; Beatriz, el día de su boda con Roberto Alessandri; Beatriz, poco después del divorcio, en un almuerzo del Club Hípico; Beatriz, en Quilmes, con Delia San Marco Porcel y Carlos Argentino; Beatriz, con el pekinés que le regaló Villegas Haedo; Beatriz, de frente y de tres cuartos, sonriendo, la mano en el mentón…
\end{quotation}

El precio que paga por esas imágenes es alto. Carlos Argentino es la encarnación viviente de todo lo que odia Borges, capaz de presumir de estrofas como esta:

\begin{verse}
A manderecha del poste rutinario \\
(viniendo, claro está, desde el Nornoroeste) \\
se aburre una osamenta ---¿Color? Blanquiceleste; \\
que da al corral de ovejas catadura de osario.
\end{verse}

Cuatro alejandrinos que condensan la insensatez de la empresa que el poetastro pretende acometer: un poema, titulado «La Tierra», que aspira a describir todo el planeta. No contento con su afán descabellado, Argentino parece empeñado en agotar todas las variaciones del mal gusto:

\begin{quotation}
[…] donde antes escribió azulado, ahora abundaba en azulino, azulenco y hasta azulillo. La palabra lechoso no era bastante fea para él; en la impetuosa descripción de un lavadero de lanas, prefería lactario, lacticinoso, lactescente, lechal…
\end{quotation}

Borges (y el lector) soporta al individuo a duras penas y todo el relato parece abocarse a un melancólico sinsentido… Hasta que de repente, como caído del cielo, aparece un Aleph.

Borges se entera de esa circunstancia cuando Argentino le informa de que sus caseros planean demoler el edificio donde vive, con ánimo de ampliar una confitería. Para el escritor, la noticia es un golpe bajo, ya que con la destrucción del inmueble se difumina un poco más la memoria de Beatriz (de nuevo la muerte térmica del Universo). El poetastro en cambio, se queja de que, para terminar su poema, le es imprescindible el edificio, ya que en un ángulo del sótano se oculta un Aleph. Y explica, de paso, que el Aleph es uno de los puntos del espacio que contiene todos los puntos, el lugar donde están, sin confundirse, todos los lugares del orbe, vistos desde todos los ángulos.

Aquí el relato se torna onírico. Un Borges incrédulo relata la rabieta del niño malcriado, convencido de que el Aleph en cuestión, además de ser de su propiedad, existe con el único propósito de ayudarle a componer su poema universal. La situación se resuelve cuando el gran escritor se decide a comprobar la patraña por sí mismo… y se encuentra con el prodigio, que nos describe con su prosa exuberante:

\begin{quotation}
Cada cosa (la luna del espejo, digamos) era infinitas cosas, porque yo claramente la veía desde todos los puntos del universo. Vi el populoso mar, vi el alba y la tarde, vi las muchedumbres de América, […] vi todos los espejos del planeta y ninguno me reflejó, […] vi tigres, émbolos, bisontes, marejadas y ejércitos, vi todas las hormigas que hay en la tierra, […] vi en un cajón del escritorio (y la letra me hizo temblar) cartas obscenas, increíbles, precisas, que Beatriz había dirigido a Carlos Argentino…
\end{quotation}

La enumeración de las maravillas que muestra el Aleph es el polo opuesto a las mamarrachadas que ha escrito el primo hermano y, ¡ay!, objeto de los desvaríos eróticos de una Beatriz a la que el punto que contiene todos los puntos despoja, de un zarpazo, del lugar en el Paraíso donde Borges y Dante la habían colocado para transformarla en una mujer de carne y hueso, cuya ausencia, por esa precisa (y obscena) carnalidad, duele todavía más. La parodia ha terminado y el recorrido a través del espejo está tintado de dolor.

Y sin embargo, el prodigio que Borges nombra con la primera letra del alfabeto hebreo ---la letra que contiene todas las demás--- admite una lectura que el propio escritor, hijo de su tiempo, no podía formular del todo pero quizás intuyó: un punto que contiene todos los puntos describe exactamente el instante infinitesimal que precede al Big Bang, la gran explosión que creó el universo hace unos trece mil ochocientos millones de años. De ahí, de ese punto, ese cero, ese lugar del orbe que contiene todos los lugares del orbe, emanan materia y antimateria (esta última desaparece enseguida, por culpa de un traidor, cuya pista ha seguido, durante toda su carrera, el que escribe estas líneas), leptones y quarks que se confinan en protones que forman núcleos de hidrógeno y después de helio y después de litio y más tarde berilio… y luego estrellas que explotan en supernovas, que crean sistemas solares donde aparece (inexplicablemente) la vida ---o al menos eso ocurre en uno de ellos---: bacterias que, casi por arte de magia, engendran eucariotas, eucariotas que engendran algas, algas que engendran nautilus, y ciempiés, medusas y tiburones, helechos y plantas carnívoras, tiranosaurios y cocodrilos, tigres de dientes de sable, caballos, unicornios y mariposas. Y en algún momento, ciegamente, aparece también una especie de monos locos, capaces de inventar la literatura para acallar el ruido en el interior de su cabeza, la poesía para mitigar su miedo a la muerte, la ciencia para tratar de entender el mundo.

Pero todo eso ocurre después. El Big Bang crea el universo pero el Aleph lo contiene en el tiempo infinitesimal que lo precede. Y ese instante infinitesimal, dicen los físicos, es una fluctuación del vacío cuántico.

Dos tremendas palabras. Vacío. Cuántico.

En el vacío clásico no hay nada. La teoría de la relatividad nos informa de que la materia teje la estructura misma del espacio-tiempo, así que el vacío, que prohíbe la materia, tampoco puede contener espacio ni tiempo. Es una ausencia absoluta.

Pero las leyes de la mecánica cuántica nos aseguran que tal vacío no existe. Aseguran que en el tejido del cosmos, incluso en la ausencia de toda materia, se forman y se destruyen, incesantemente, pares de partículas y antipartículas, en un baile invisible, oculto al observador por el implacable principio de incertidumbre. Y ese fenómeno puede fluctuar, generar una burbuja, provocar un Big Bang, quizás con una probabilidad tan pequeña que sería imposible que ocurriera, a no ser que se disponga de toda la eternidad para ello.

Y así, eventualmente, quizás en un inconcebible Multiverso, aparecen burbujas iniciadas por esa rebeldía del vacío cuántico que forman un Aleph, que a su vez engendra el orbe.

Hacia el final de su magistral relato, Borges nos informa que el Aleph que existía en el sótano del poetastro Carlos Argentino era, posiblemente, un falso Aleph. Es fácil creer que está en lo cierto y aventurar que el auténtico Aleph desapareció hace miles de millones de años, dando lugar a todas las cosas.

O quizás no es así. Quizás el vacío cuántico sigue fluctuando y continuamente se forman otros universos que se desprenden del nuestro como burbujas de jabón sopladas por el lapicero hueco de un Dios travieso. Quizás todos hemos visto un Aleph sin saberlo, porque esa visión, demasiado sobrecogedora, se ha borrado de nuestra mente, como se han borrado ya, para siempre, los rasgos de Beatriz Viterbo.

\end{document}
